\ncertchapter{14}{Torque, Angular Momentum and Moment of Inertia}\chlabel{ch:rotation}

Push a door at its handle and it swings open. Push it just as hard at the hinge and nothing happens. Force,
which has governed everything so far, is plainly not the whole story once a body can rotate: \emph{where} the
force is applied matters as much as how large it is. This chapter builds the mechanics of rotation, and it does
so by constructing, one by one, the rotational partner of every quantity in translational mechanics.

Three new quantities carry the chapter, and each earns its place by answering a different practical question.
\emph{Torque} measures the turning effect of a force, and explains why a long spanner loosens a stiff nut that
a short one cannot. \emph{Moment of inertia} measures a body's resistance to angular acceleration, and explains
why a figure skater spins faster with her arms drawn in. \emph{Angular momentum} is the rotational partner of
linear momentum and is conserved when no external torque acts; it is the reason a gyroscope holds its orientation and a collapsing
star core spins up to tens of revolutions per second, and the way a torque changes it explains why a spinning top precesses instead of falling over. (It is not,
despite a popular myth, what keeps a moving bicycle upright: steering geometry matters far more.)

The central equation is $\tau=I\alpha$, the rotational form of Newton's second law, and the chapter is careful
about where it holds. It also derives the standard moments of inertia rather than quoting them, and proves
the parallel- and perpendicular-axis theorems that turn a short table into a moment of inertia about any axis.

\begin{remarkbox}{Notation in this chapter}
Rotational mechanics reuses a few letters, so fix them now.
\begin{itemize}
\item $\vL$ (and its magnitude $L$) is \textbf{angular momentum} throughout this chapter, as is
      universal usage.
\item Because of that, the \textbf{length of a rod} is written $\ell$ here, not $L$ as in the earlier
      chapters. Where you see $M\ell^{2}$ it is a length; where you see $L=I\omega$ it is angular
      momentum.
\item $I$ is always a moment of inertia, with $I_{\text{cm}}$ reserved for the axis through the centre
      of mass.
\item $\alpha$ remains \textbf{angular acceleration}, never an angle, exactly as in \chsref{ch:circkin}{ch:circdyn}.
\item $r_\perp$ is a perpendicular distance from the axis (or chosen point): to a force's line of action
      for torque, and to a particle's line of motion for angular momentum.
\item $\theta$ does two jobs, kept apart by context: the angle \emph{between two vectors} in a cross
      product or torque, and the \emph{angular position} of a rotating body in kinematics.
\end{itemize}
\end{remarkbox}

\section{The Cross Product}
Rotation forces a new piece of vector algebra on us. Up to now the only product of two vectors we needed was the dot
product, which answers ``how much of one vector lies \emph{along} another''. But turning is about
sideways-ness: a force pushing straight at a hinge turns nothing, however hard it pushes, while the same
force applied across the hinge turns it easily. We therefore need a product that measures how much of
one vector lies \emph{across} another, and that returns a direction, since a rotation has an axis. That
product is the cross product, and everything in this chapter is built from it.

\begin{defbox}{Vector (cross) product}
$\A\ex\B=AB\sin\theta\,\uhat{n}$, where $A$ and $B$ are the magnitudes, $\theta$ ($0\le\theta\le\pi$) is
the angle between the vectors, and $\uhat{n}$ is a unit vector. The product is a \emph{vector} of magnitude $AB\sin\theta$, perpendicular to both $\A$
and $\B$, pointing along the direction your right thumb gives when the fingers curl from $\A$ to $\B$. It
is anti-commutative ($\A\ex\B=-\B\ex\A$) and vanishes when the vectors are parallel. In components it is
the determinant below, with $\ihat,\jhat,\khat$ the unit vectors along the axes (\chref{ch:vectors}):
\[
\A\ex\B=\begin{vmatrix}\ihat & \jhat & \khat\\ A_x & A_y & A_z\\ B_x & B_y & B_z\end{vmatrix},
\qquad\text{with}\quad \ihat\ex\jhat=\khat,\ \ \jhat\ex\khat=\ihat,\ \ \khat\ex\ihat=\jhat.
\]
\end{defbox}

\begin{figure}[htbp]
\centering
\tikzfig{cross_product}
\figcaption{fig:cross}{Where the dot product $\A\cdot\B=AB\cos\theta$ asks how much of $\A$ lies \emph{along} $\B$, the cross product asks how much lies \emph{across}. Left: its magnitude is the area of the parallelogram the two vectors span, base $A$ times height $B\sin\theta$, largest when they are perpendicular and zero when they are parallel. Right: its direction is perpendicular to the plane of $\A$ and $\B$, given by the right thumb when the fingers curl from $\A$ to $\B$; reversing the order reverses it, $\B\ex\A=-\A\ex\B$.}
\end{figure}

\why{\textbf{Why the result is a vector, and a strange one.} A rotation needs an axis, and an axis is
a direction, so the natural way to encode ``rotation about this line, this fast'' is a vector \emph{along}
that line. The right-hand rule that fixes which way it points is only a convention: a left-hand rule
would work just as well if applied consistently. Torque and angular momentum (met below) are vectors of this kind, as is the angular velocity $\vomega$
of \chref{ch:circkin}, which points along the axis. This also ties off a loose end from \chref{ch:vectors}, where we saw that finite rotations are \emph{not} vectors
because they do not commute (turning a book $90^\circ$ about two axes in the two orders gives different
results). The resolution: \emph{infinitesimal} rotations \emph{do} commute to first order, and $\vomega$
is built from those, which is why $\vomega$ is a legal vector even though ``rotate by $90^\circ$'' is
not.}

\section{Torque}
A force makes a body accelerate. What makes a body \emph{spin}? Everyday experience says it is not the
force alone: pushing a door at its handle swings it open, pushing with identical force right beside the
hinge does almost nothing, and pushing directly towards the hinge does nothing at all. So the turning
effect depends on \emph{where} and \emph{in what direction} the force is applied, not merely on its
size. The quantity that captures all three facts at once is torque. First, though, we must say what kind
of body is doing the turning.

\begin{defbox}{Rigid body and fixed-axis rotation}
A \textbf{rigid body} is one whose particles keep fixed distances from one another, however it moves. It
is an idealisation (every real body bends a little), but a good one for wheels, rods and doors. When a
rigid body turns about a \textbf{fixed axis}, every particle moves in a circle centred on that axis, and
all of them turn through the \emph{same} angle in the same time. So one angle $\theta$, one angular
velocity $\omega$ and one angular acceleration $\alpha$ (\chref{ch:circkin}) describe the motion of the
whole body, and that is what makes rotation tractable. Most of this chapter is about fixed-axis rotation;
\chref{ch:rolling} lets the axis move.
\end{defbox}

\begin{motivbox}{Why torque matters}
Force alone answers ``will it move?'' Torque answers ``will it \emph{turn}?'' It is why spanners, door
handles and bicycle cranks are long, why an engine is rated by its torque in newton-metres, and why a
crane or a loaded truck tips over when the torque of its load about the tipping edge exceeds the
restoring torque of its own weight.
\end{motivbox}

With a rigid body in hand, the turning effect of a force can be defined. It must grow with the force,
grow with the distance from the turning point, and vanish when the force points straight at that point:
the cross product of the previous section does exactly this.

\begin{defbox}{Torque}
The torque of a force $\vF$ about a point is $\vtau=\vr\ex\vF$, where $\vr$ runs from that point to where the force acts, with magnitude
\[
\tau=rF\sin\theta=F\,r_\perp,
\]
where $\theta$ is the angle between $\vr$ and $\vF$, and $r_\perp=r\sin\theta$ is the \textbf{moment arm}: the perpendicular distance from the point to
the force's \textbf{line of action} (the infinite line through the point of application, along $\vF$). Units: newton-metres (never joules, though the dimensions match).
\end{defbox}

\textbf{About a point, or about an axis?} Torque is defined about a \emph{point}, and it changes if the
point changes: the same force has a large torque about a distant point and none about a point on its line
of action, so always name the point before calculating. For a body turning about an axis, what matters is
the component of $\vtau$ (taken about any point on the axis) along the axis. In every problem in this
chapter the forces lie in the plane perpendicular to the axis, and then that component is simply
$\pm F\,r_\perp$, with $r_\perp$ measured from the axis (Figure~\ref{fig:torque} shows how it shrinks as the push turns towards the hinge): by the right-hand rule, an anticlockwise torque (seen
from the side the axis points towards) is out of the page, and we take it positive. A force parallel to
the axis, or one whose line of action passes through it, has no turning effect about that axis.

\begin{figure}[htbp]
\centering
\tikzfig{torque_door}
\figcaption{fig:torque}{A door seen from above, its hinge axis out of the page. The torque is $F\,r_\perp$, where the moment arm $r_\perp$ is the perpendicular distance from the axis to the force's line of action. Pushed square on at the free edge, $r_\perp=r$ and the torque is largest; at an angle $\theta$ to the door, $r_\perp=r\sin\theta$; pushed along the door ($\theta=0$) or at the hinge ($r=0$), the line of action passes through the axis and the torque vanishes. This is why handles sit at the far edge and long spanners undo tight bolts.}
\end{figure}

\begin{trapbox}{torque and energy share dimensions but nothing else}
Both torque and energy have dimensions $\mathsf M\mathsf L^2\mathsf T^{-2}$, yet torque is a
\emph{vector} built from a cross product and energy is a \emph{scalar} built from a dot product. They may
never be added, and dimensional analysis (\chref{ch:measurement}) cannot tell them apart.
\end{trapbox}

\inwords{The moment arm explains the seesaw with no new physics. Balance means the two torques about the
pivot are equal, $m_1 g\,r_1=m_2 g\,r_2$, where $r_1$ and $r_2$ are the two sitters' distances from the pivot: it is mass \emph{times distance} that must match, not mass.
So a heavy adult sits close to the pivot and a light child far out.}

A particle is in equilibrium when the net force on it is zero (\chref{ch:newton}). A rigid body needs
more, and the simplest case shows why. A \textbf{couple} is two equal and opposite forces whose lines of
action do not coincide, separated by a perpendicular distance $d$. Their net force is zero, so the centre of
mass does not accelerate, yet the body spins up: turning a steering wheel with both hands, or a tap with
finger and thumb, applies a couple. Taking moments about any point $O$ on the line of one force, the
other has arm $d$, so $\tau=Fd$; moving $O$ elsewhere adds as much torque from one force as it removes
from the other, so a couple has the same torque $Fd$ about \emph{every} point. Zero net force is
therefore not enough for a rigid body to stay at rest.

\begin{keybox}{Equilibrium of a rigid body and the principle of moments}
A rigid body at rest stays at rest only if \emph{both}
\[
\textstyle\sum\vF_{\text{ext}}=\vc{0}\quad\text{(no translation)}\qquad\text{and}\qquad
\sum\vtau_{\text{ext}}=\vc{0}\quad\text{(no rotation)}.
\]
The second is the \textbf{principle of moments}: about any point, the total clockwise torque equals the
total anticlockwise torque. When the net force is zero, the net torque is the same about every point
(shifting the point by $\vc{b}$ turns every $\vr$ into $\vr-\vc{b}$, which changes the total torque by $-\vc{b}\ex\sum\vF$), so the point may be
chosen to make the arithmetic easiest, usually through an unknown force, which then drops out. A body's
weight acts at its centre of mass (\chref{ch:com}). For a lever this is
$F_{\text{load}}\,d_{\text{load}}=F_{\text{effort}}\,d_{\text{effort}}$, giving a mechanical advantage
$d_{\text{effort}}/d_{\text{load}}$: a small effort raises a large load, at the cost of distance moved
(\chref{ch:connected}'s movable pulley in rotational dress).
\end{keybox}

\begin{trapbox}{the freedom to choose any point is an equilibrium privilege}
For a body that is \emph{accelerating}, the net force is not zero, and torques about different points
differ. Moments may then be taken only about a fixed axis or about the centre of mass (see the scope of
$\tau=I\alpha$ below), which is why \chref{ch:circdyn} took moments about $G$ for the cornering cyclist.
\end{trapbox}

\section{Moment of Inertia}

Torque tells us what makes a body's spin \emph{change}; the next question is how strongly the body
\emph{resists} that change. For straight-line motion the answer was mass. For rotation it is a new
quantity: the moment of inertia, and, unlike mass, it depends on where the axis is.

\begin{motivbox}{Why moment of inertia matters}
Moment of inertia depends not only on how much mass there is but on \emph{how far from the axis} it sits. That one fact
is why a flywheel is built heavy at the rim, why a tightrope walker's long pole slows a topple enough to
correct, and why a diver tucks to spin faster.
\end{motivbox}

Why distance should matter, and squared, is easy to see. To give a particle of mass $m$ at distance $r$
from the axis an angular acceleration $\alpha$, its tangential acceleration $\alpha r$ needs a force
$m\alpha r$, whose torque about the axis is $m\alpha r\cdot r=mr^2\alpha$: one factor of $r$ from the
acceleration, one from the moment arm.

\begin{defbox}{Moment of inertia}
$I=\sum_i m_i r_i^2$ for particles of mass $m_i$, or $\int r^2\,dm$ for a continuous body cut into small
mass elements $dm$ (\chref{ch:com}), where $r$ is the perpendicular distance of each from the axis. It is a
scalar, the rotational analogue of mass (the resistance to angular acceleration), but it is a
\emph{number per axis}, not a fixed property of the body. Units: \si{kg.m^2}.
\end{defbox}

The square makes distance dominate: Figure~\ref{fig:moir2} places the same mass at $r$, $2r$ and $3r$ and compares the contributions.

\begin{figure}[htbp]
\centering
\tikzfig{moi_r2}
\figcaption{fig:moir2}{Because $I$ weights each mass element by $r^2$, distance dominates: mass at three
times the radius counts \emph{nine} times as much. This is why flywheels are built heavy at the rim, why
a hollow cylinder resists spinning more than a solid one of equal mass, and why a tightrope walker's long
pole, with its mass far from the rope, makes toppling slow enough to correct.}
\end{figure}

\begin{derivbox}{A thin rod about its centre}
Take a uniform rod of mass $M$, length $\ell$, spun about a perpendicular axis through its centre. Its linear
density is $\lambda=M/\ell$, and a small element $dx$ at distance $x$ from the axis (Figure~\ref{fig:moidrod}) has mass $dm=\lambda\,dx$
and contributes $x^2\,dm$. Sum (integrate) over the whole rod, from $-\ell/2$ to $+\ell/2$:
\[
I=\int_{-\ell/2}^{\ell/2} x^2\lambda\,dx=\lambda\left[\frac{x^3}{3}\right]_{-\ell/2}^{\ell/2}
=\lambda\cdot\frac{\ell^3}{12}=\frac{M}{\ell}\cdot\frac{\ell^3}{12}=\frac{1}{12}M\ell^2.
\]
\end{derivbox}

\begin{figure}[htbp]
\centering
\tikzfig{moid_rod}
\figcaption{fig:moidrod}{The slice for the rod: an element $dx$ at signed distance $x$ from the central
axis, mass $dm=\lambda\,dx$. Integrating $x^2\,dm$ from $-\ell/2$ to $\ell/2$ gives $\tfrac{1}{12}M\ell^2$; moving
the axis to one end (limits $0$ to $\ell$) gives $\tfrac{1}{3} M\ell^2$.}
\end{figure}

The same slicing gives every standard shape. The results are collected here so they sit in one place;
each is derived in the subsection that follows.

\begin{tablebox}{Standard moments of inertia ($M$ mass, $R$ radius; $I=cMR^2$, or $cM\ell^2$ for a rod)}
\centering
\begin{tabular}{@{}ll@{\hspace{2em}}ll@{}}
\toprule
Ring (central axis) & $MR^2$ & Solid sphere & $\tfrac{2}{5} MR^2$\\
Ring (diameter) & $\tfrac{1}{2} MR^2$ & Hollow sphere (shell) & $\tfrac{2}{3} MR^2$\\
Disc / solid cylinder & $\tfrac{1}{2} MR^2$ & Rod (through centre) & $\tfrac{1}{12}M\ell^2$\\
Rod (through end) & $\tfrac{1}{3} M\ell^2$ & &\\
\bottomrule
\end{tabular}
\end{tablebox}

Read the number $c$, the \textbf{shape factor}, as a spread score: the larger it is, the farther out the mass sits. A ring has $c=1$
(all its mass at the maximum radius $R$); a solid sphere has $c=\tfrac{2}{5}$ (mass hugging the centre). The
ordering $\tfrac{2}{5}<\tfrac{1}{2}<\tfrac{2}{3}<1$ is worth knowing cold; \chref{ch:rolling} shows that it decides a race.
The \textbf{radius of gyration} $k$ says the same thing as a length: it is the radius at which a point mass $M$
would have the same $I$, so $Mk^2=I$ and $k=\sqrt{I/M}=\sqrt{c}\,R$. A disc
($k=R/\sqrt2\approx0.71R$) resists turning exactly as a ring of the same mass and radius $0.71R$ would.
Notice also that $c$ does not involve $M$ or $R$: two discs of different sizes and masses share $c=\tfrac12$.

\subsection{Deriving the standard moments of inertia}

Each follows the same recipe as the centres of mass in \chref{ch:com}, but now we integrate $r^2\,dm$: the perpendicular
distance from the axis, \emph{squared}. Choose slices on which $r$ is constant and the integral collapses
to one of \chref{ch:com}'s three mass elements: $dm=\lambda\,dx$ along a line, $\sigma\,dA$ over a surface or $\rho\,dV$ through a volume, with $\lambda$, $\sigma$, $\rho$ the mass per unit length, area and volume.

\begin{derivbox}{Ring about its central axis, $I=MR^2$ (no integral)}
Every mass element of a ring (or thin hoop) sits at the \emph{same} distance $R$ from the central axis
(Figure~\ref{fig:moidring}, left).
So $r=R$ is constant and comes straight out of the sum: $I=\sum m_i R^2=R^2\sum m_i=MR^2$. This is the
$c=1$ benchmark: the most spread-out mass distribution possible for a given radius, and the reason a
hoop is the hardest shape to spin up.
\end{derivbox}

\begin{figure}[htbp]
\centering
\tikzfig{moid_ring}
\figcaption{fig:moidring}{The ring two ways. About the \emph{central} axis (left) every element sits at
the same distance $R$, so $I=MR^2$ with no integral. About a \emph{diameter} (right) the perpendicular-axis
theorem splits $I_z=MR^2$ into two equal in-plane parts, giving $\tfrac{1}{2} MR^2$.}
\end{figure}

\begin{derivbox}{Disc / solid cylinder about its axis, $I=\tfrac{1}{2} MR^2$}
A disc is a nest of rings. Slice it into thin rings of radius $r$ and width $dr$ (Figure~\ref{fig:moiddisc}); each has
$dm=\sigma\,2\pi r\,dr$ (with $\sigma=M/\pi R^2$) and, being a ring, contributes $dI=r^2\,dm$:
\[
I=\int_0^R r^2\,\sigma\,2\pi r\,dr=2\pi\sigma\int_0^R r^3\,dr=2\pi\sigma\,\frac{R^4}{4}
=2\pi\cdot\frac{M}{\pi R^2}\cdot\frac{R^4}{4}=\frac{1}{2}MR^2.
\]
A solid cylinder is just a stack of identical discs, so its axial moment of inertia is the same,
$\tfrac{1}{2} MR^2$: the length does not enter. The $\tfrac{1}{2}$ (against the ring's $1$) reflects that a
disc's mass is spread inward from the rim, not all at $R$.
\end{derivbox}

\begin{figure}[htbp]
\centering
\tikzfig{moid_disc}
\figcaption{fig:moiddisc}{The slice for the disc: a thin ring of radius $r$ and width $dr$. Every point of the ring is at the same distance $r$ from the axis, so its contribution is $r^2\,dm$, and summing the rings from $0$ to $R$ gives $\tfrac{1}{2} MR^2$.}
\end{figure}

\begin{derivbox}{Solid sphere about a diameter, $I=\tfrac{2}{5} MR^2$}
Slice the sphere into thin discs perpendicular to the chosen axis (Figure~\ref{fig:moidsphere}). A disc at height $z$ has radius
$a=\sqrt{R^2-z^2}$ and thickness $dz$, so its mass is $dm=\rho\,\pi a^2\,dz$ (with
$\rho=M/\tfrac{4}{3}\pi R^3$). A disc's own moment of inertia about its central axis is $\tfrac{1}{2}(dm)a^2$
(from the previous box), and here that disc's axis \emph{is} the sphere's axis, so
\[
I=\int_{-R}^{R}\tfrac{1}{2} a^2\,dm=\tfrac{1}{2}\rho\pi\int_{-R}^{R}(R^2-z^2)^2\,dz.
\]
Expanding $(R^2-z^2)^2=R^4-2R^2z^2+z^4$ and integrating from $-R$ to $R$ gives
$2R^5-\tfrac{4R^5}{3}+\tfrac{2R^5}{5}=\tfrac{16}{15}R^5$. Then
$I=\tfrac{1}{2}\rho\pi\cdot\tfrac{16}{15}R^5$, and substituting $\rho=\dfrac{3M}{4\pi R^3}$ collapses this to
$\boxed{I=\tfrac{2}{5} MR^2}$.
\end{derivbox}

\Needspace*{20\baselineskip}
\begin{figure}[htbp]
\centering
\tikzfig{moid_sphere}
\figcaption{fig:moidsphere}{The slice for the solid sphere: a disc of thickness $dz$ at height $z$, of radius $a=\sqrt{R^2-z^2}$. Each disc contributes $\tfrac{1}{2} a^2\,dm$ (its own moment of inertia about the same axis), and stacking them from $-R$ to $R$ gives $\tfrac{2}{5} MR^2$.}
\end{figure}

\begin{derivbox}{Hollow sphere (thin shell) about a diameter, $I=\tfrac{2}{3} MR^2$}
Slice the shell into rings by the polar angle $\theta$ from the axis (Figure~\ref{fig:moidhollow}). A ring at angle $\theta$ has radius
$r=R\sin\theta$ (its distance from the axis), slant width $R\,d\theta$, and circumference
$2\pi R\sin\theta$, so its mass is $dm=\sigma\,(2\pi R\sin\theta)(R\,d\theta)$ with $\sigma=M/4\pi R^2$.
Being a ring, it contributes $dI=r^2\,dm=(R\sin\theta)^2\,dm$:
\[
I=\int_0^\pi (R\sin\theta)^2\,\sigma\,2\pi R^2\sin\theta\,d\theta
=2\pi\sigma R^4\int_0^\pi \sin^3\theta\,d\theta=2\pi\sigma R^4\cdot\frac{4}{3},
\]
where the integral is done by writing $\sin^3\theta=\sin\theta\,(1-\cos^2\theta)$ and substituting $u=\cos\theta$.
Substituting $\sigma=\dfrac{M}{4\pi R^2}$ gives $\boxed{I=\tfrac{2}{3} MR^2}$. It exceeds the solid sphere's
$\tfrac{2}{5}$ because a shell keeps \emph{all} its mass out at the surface, none near the centre: the same
reason a hoop beats a disc.
\end{derivbox}

\begin{figure}[htbp]
\centering
\tikzfig{moid_hollowsphere}
\figcaption{fig:moidhollow}{The slice for the spherical shell: a ring traced by the polar angle $\theta$,
at distance $r=R\sin\theta$ from the axis. Being a ring, it contributes $r^2\,dm$, and integrating over
$\theta$ (via $\int_0^\pi\sin^3\theta\,d\theta=\tfrac{4}{3}$) gives $\tfrac{2}{3} MR^2$.}
\end{figure}

\begin{derivbox}{Rod about one end, $I=\tfrac{1}{3} M\ell^2$ (two ways)}
\textbf{By integration:} with the axis at one end, $I=\int_0^{\ell} x^2\lambda\,dx=\lambda\dfrac{\ell^3}{3}=\tfrac{1}{3} M\ell^2$
(the same setup as the centre case, but limits $0$ to $\ell$). \textbf{By the parallel-axis theorem} (next
section), it is instant: shift from the centre by $d=\ell/2$,
$I_{\text{end}}=I_{\text{cm}}+M\left(\tfrac{\ell}{2}\right)^2=\tfrac{1}{12}M\ell^2+\tfrac{1}{4} M\ell^2=\tfrac{1}{3} M\ell^2$.
The two agreeing is a good check that the theorem does what it claims.
\end{derivbox}

\section{The Two Theorems}
Integration is laborious, and for most examination problems unnecessary. The two theorems below let you take a standard result you
already know and transport it to a new axis in a single line. One handles a shift to a parallel axis and
works for any body at all; the other relates three mutually perpendicular axes but works only for flat
bodies. Figure~\ref{fig:theorems} sets the two side by side.

\begin{derivbox}{The parallel axis theorem}
$I=I_{\text{cm}}+Md^2$, where $d$ is the distance between the axis and a parallel axis through the centre of mass (CM).
\stepn{1} Put the origin at the CM with the $z$-axis along the CM axis, and place the new axis parallel to it, a distance $d$ away along $x$. A mass element at $(x,y,z)$ is at distance-squared $x^2+y^2$
from the CM axis (its $z$-coordinate is irrelevant), and $(x-d)^2+y^2$ from the new axis. So
\[
I=\int\!\big[(x-d)^2+y^2\big]dm=\underbrace{\int(x^2+y^2)\,dm}_{I_{\text{cm}}}
-2d\underbrace{\int x\,dm}_{=\,Mx_{\text{cm}}=0}+d^2\underbrace{\int dm}_{M}.
\]
\stepn{2} \textbf{The middle term vanishes} precisely because the origin is the CM, so
$\int x\,dm=Mx_{\text{cm}}=0$. That is the whole trick, and it is why the theorem needs the CM axis
specifically. Hence $I=I_{\text{cm}}+Md^2$. To move between two axes \emph{neither} of which passes
through the CM, go via the CM: $I_2=I_1-Md_1^2+Md_2^2$, not $I_1+M(d_2-d_1)^2$.
\stepn{3} \textbf{A free gift.} Since $Md^2\ge0$ always, $I$ is \emph{smallest} about an axis through the
CM.
\end{derivbox}

\begin{derivbox}{The perpendicular axis theorem (laminae only)}
For a flat lamina in the $xy$-plane, $I_z=I_x+I_y$. A point in the plane is at distance $|y|$ from the
$x$-axis, $|x|$ from the $y$-axis, and $r=\sqrt{x^2+y^2}$ from the $z$-axis. So
\[
I_z=\int r^2\,dm=\int(x^2+y^2)\,dm=\underbrace{\int y^2\,dm}_{I_x}+\underbrace{\int x^2\,dm}_{I_y}.
\]
The proof needs $r^2=x^2+y^2$, which requires $z=0$, \emph{all} the mass in one plane. Applying it to
a solid sphere or cylinder is a standard and fatal error. The three axes must also meet at one point,
with $x$ and $y$ in the plane and $z$ perpendicular to it; that point need not be the CM.
\end{derivbox}

\begin{figure}[htbp]
\centering
\tikzfig{moi_theorems}
\figcaption{fig:theorems}{The two theorems. \emph{Parallel axis} shifts $I$ from the CM to any parallel
axis, its cross-term vanishing because the origin is the CM. \emph{Perpendicular axis} combines two
in-plane axes into the one perpendicular to the lamina, and needs every scrap of mass in the plane, so
it fails for solids.}
\end{figure}

\begin{methodbox}{Using the theorems together}
For a \emph{ring about its diameter}: by symmetry $I_x=I_y$, and the perpendicular-axis theorem gives
$I_z=MR^2=I_x+I_y=2I_x$, so $I_{\text{diameter}}=\tfrac{1}{2} MR^2$, with no integral at all. For
composite bodies, parallel-axis each piece to the common axis and add (cut-outs subtract).
\end{methodbox}

\section{\texorpdfstring{Rotational Dynamics: $\tau=I\alpha$}{Rotational Dynamics: torque = I alpha}}

We now have every ingredient: torque measures the turning effort, and moment of inertia measures the
resistance to being turned. Newton's second law for rotation connects them, and it is the
most used equation in this chapter. We recall the rotational kinematics first, then derive the
dynamics.

\begin{keybox}{Angular kinematics at constant $\alpha$ (from \chref{ch:circkin})}
A rotating body is described by its angular position $\theta$, angular velocity $\omega=d\theta/dt$, and
angular acceleration $\alpha=d\omega/dt$: the exact rotational partners of $x$, $v$, $a$. When
$\alpha$ is \emph{constant}, they obey the same three equations as straight-line motion, with the letters
swapped:
\[
\omega=\omega_0+\alpha t,\qquad
\theta=\omega_0 t+\tfrac{1}{2}\alpha t^2,\qquad
\omega^2=\omega_0^2+2\alpha\theta.
\]
As a vector, $\vomega$ points along the axis, the way the right thumb points when the fingers curl
with the rotation; $\vc{\alpha}=d\vomega/dt$. A point at radius $r$ from the axis has linear speed $v=\omega r$ and tangential acceleration $a=\alpha r$,
which is how the rotation of the whole body connects to the motion of its parts.
\end{keybox}

Before deriving $\tau=I\alpha$ it is worth seeing how completely rotation mirrors translation,
because that mapping turns a subject which looks new into one you already know. The angular-momentum rows are previewed
here and derived in \S\ref{sec:angmom14}.

\begin{tablebox}{The translation--rotation dictionary}
\centering
\begin{tabular}{@{}ll@{\hspace{3em}}ll@{}}
\toprule
\multicolumn{2}{c}{\textbf{Translation}} & \multicolumn{2}{c}{\textbf{Rotation}}\\
position & $x$ & angle & $\theta$\\
velocity & $v$ & angular velocity & $\omega$\\
mass & $m$ & moment of inertia & $I$\\
force & $\vF$ & torque & $\vtau=\vr\ex\vF$\\
momentum & $\vp=m\vv$ & angular momentum & $\vL=I\vomega$\\
$\vF=d\vp/dt$ & & $\vtau=d\vL/dt$ &\\
$K=\tfrac{1}{2}mv^2$ & & $K=\tfrac{1}{2}I\omega^2$ &\\
work & $W=\int\vF\cdot d\vr$ & work & $W=\int\tau\,d\theta$\\
power & $P=\vF\cdot\vv$ & power & $P=\tau\,\omega$\\
\bottomrule
\end{tabular}
\end{tablebox}

The one genuinely new object is $I$: unlike mass, it depends on the axis. Three entries of
the dictionary, $\tau=I\alpha$, $K=\tfrac12I\omega^2$ and $W=\int\tau\,d\theta$, are worth deriving, since
they are the rotational engines of the whole subject.

\begin{derivbox}{Why $\tau=I\alpha$ and $K_{\text{rot}}=\tfrac{1}{2} I\omega^2$}
\small
Both follow by summing the particle laws over a rigid body turning about a fixed axis, using $v_i=\omega r_i$ and $a_i=\alpha r_i$
(every piece shares the same $\omega,\alpha$; $r_i$ is the distance of piece $i$ from the axis, and $\tau$ means the torque component along the axis).
\stepn{$\tau=I\alpha$} The tangential force on piece $m_i$ is $m_i a_i=m_i\alpha r_i$ (the radial part of
the force points at the axis and has no torque about it), and its torque
about the axis is that force times its arm $r_i$: $\tau_i=m_i\alpha r_i\cdot r_i=m_i r_i^2\,\alpha$.
Summing (the forces between pieces come in equal and opposite pairs along the line joining them, so
their torques cancel in pairs and only the external torque survives), $\tau=\sum_i m_i r_i^2\,\alpha=I\alpha$: the exact analogue of $F=ma$, with $I$ in the role of
mass.
\stepn{$K=\tfrac{1}{2} I\omega^2$} Piece $m_i$ moves at speed $v_i=\omega r_i$, so its kinetic energy is
$\tfrac{1}{2} m_i v_i^2=\tfrac{1}{2} m_i r_i^2\,\omega^2$. Summing,
$K=\tfrac{1}{2}\big(\sum_i m_i r_i^2\big)\omega^2=\tfrac{1}{2} I\omega^2$: the analogue of $\tfrac{1}{2} mv^2$.
\stepn{$W=\int\tau\,d\theta$} A force with tangential component $F_t$ at radius $r$ moves its point of
application through an arc $ds=r\,d\theta$ as the body turns by $d\theta$ (a radial component does no
work, being perpendicular to the motion). So $dW=F_t\,r\,d\theta=\tau\,d\theta$, and dividing by $dt$
gives the power $P=\tau\omega$. Since $\tau\,d\theta=I\alpha\,d\theta=I\omega\,d\omega$, the total work is
$\int\tau\,d\theta=\Delta\big(\tfrac{1}{2} I\omega^2\big)$: the work--energy theorem of \chref{ch:energy}
in rotational form.
\end{derivbox}

\begin{remarkbox}{Where $\tau=I\alpha$ holds, and where it does not}
The derivation assumed a \emph{rigid} body (so $I$ is constant) turning about a \emph{fixed} axis, and it
relied on internal forces acting along the lines joining particles. Three consequences.
(1) It also holds about an axis through the \emph{centre of mass}, even when the CM accelerates, provided
the axis keeps its direction; \chref{ch:rolling} uses exactly this, $\tau_{\text{cm}}=I_{\text{cm}}\alpha$. About
any other moving point it can fail.
(2) When $I$ changes, as for a skater drawing in her arms, $\tau=I\alpha$ is wrong; the general law is
$\tau=dL/dt$ (next section).
(3) A body that rolls has kinetic energy of translation as well, $K=\tfrac12Mv_{\text{cm}}^2+\tfrac12I_{\text{cm}}\omega^2$
(\chref{ch:rolling}); $\tfrac12I\omega^2$ alone is the whole story only about a fixed axis.
\end{remarkbox}

\section{Angular Momentum}\label{sec:angmom14}

The law $\tau=I\alpha$ has two blind spots. It cannot handle a body whose $I$ changes, such as a skater
drawing in her arms, and it says nothing about a lone particle that is not part of a rigid body, such as a
planet. Both are covered by the rotational partner of momentum, just as $\vF=d\vp/dt$ is more general than
$\vF=m\va$. Its conservation needs only that no external \emph{torque} acts, and a force aimed through the
axis exerts none, however large; so it lets you predict the end of a spinning process without following
the middle.

\begin{motivbox}{Why angular momentum matters}
\begin{itemize}
\item \textbf{Spacecraft orientation.} With nothing to push against, a space telescope turns itself
by spinning internal \emph{reaction wheels} one way, so that the craft rotates the other and the total
angular momentum stays fixed.
\item \textbf{Why helicopters need a tail rotor.} The engine exerts a torque on the main rotor, so the
rotor exerts an equal and opposite torque on the body, which would spin the other way. The tail rotor
pushes sideways on the air to supply an opposing external torque; lose it and the cabin whirls.
\end{itemize}
\end{motivbox}

\begin{defbox}{Angular momentum}
For a particle with linear momentum $\vp=m\vv$ at position $\vr$ from a chosen point,
$\vL=\vr\ex\vp=m(\vr\ex\vv)$, a vector of magnitude $L=mvr_\perp$. Units: \si{kg.m^2/s} (equivalently
\si{J.s}). For a rigid body rotating about a
fixed axis, the component of angular momentum along that axis is $L=I\omega$; the full vector $\vL=I\vomega$
holds when the axis is an axis of symmetry, which covers every case in this book. The rotational form of
Newton's second law (derived below) is
\[
\vtau_{\text{ext}}=\frac{d\vL}{dt},\qquad\text{so}\qquad \vtau_{\text{ext}}=0\ \Rightarrow\ \vL=\text{constant}.
\]
\end{defbox}

In words, $L=mvr_\perp$ is momentum times its ``miss distance'' $r_\perp$ from the chosen point. A car
driving straight past you has angular momentum about you, although it never rotates: angular momentum
measures circulation about a point, and a body need not be turning to have it.

\begin{derivbox}{Why $L=I\omega$ for a rigid body}
Chop a body turning about a fixed axis into pieces $m_i$, each at perpendicular distance $r_i$ from the
axis. Each piece moves in a circle about the axis, so its velocity is perpendicular to $r_i$ and its
angular momentum along the axis is $m_iv_ir_i$. The body is rigid, so every piece shares the same $\omega$
and $v_i=\omega r_i$:
\[
L=\sum_i m_i r_i^2\,\omega=\Big(\sum_i m_i r_i^2\Big)\omega=I\omega,
\]
the mirror of $p=mv$.
\end{derivbox}

\begin{derivbox}{Why $\vtau_{\text{ext}}=d\vL/dt$}
For one particle, differentiate $\vL=\vr\ex\vp$ by the product rule:
$\dfrac{d\vL}{dt}=\vv\ex\vp+\vr\ex\dfrac{d\vp}{dt}=\vc{0}+\vr\ex\vF=\vtau$, since $\vv\ex m\vv=\vc{0}$
(parallel vectors) and $d\vp/dt=\vF$ (\chref{ch:newton}). For a system, add this over all particles. The
internal forces come in third-law pairs acting along the line joining the two particles, so each pair has
the same moment arm and opposite senses, and their torques cancel; only external torques survive. Both
$\vtau$ and $\vL$ must be taken about the \emph{same} fixed point (or about the CM). For a rigid body with
constant $I$, $dL/dt=I\,d\omega/dt=I\alpha$, which recovers $\tau=I\alpha$ as a special case.
\end{derivbox}

\subsection{The direction of angular momentum, and why tops stay up}

So far we have used only the size of $\vL$; its direction carries just as much physics. For a spinning body, $\vL$ points \emph{along the axis of rotation}, in the sense given by the right hand:
curl your fingers the way the body turns, and your thumb points along $\vL$. No particle moves in that
direction (every particle circles \emph{around} the axis), but it is the one direction the spinning
singles out.

\begin{remarkbox}{Why a spinning body resists tipping}
To \emph{change} $\vL$ you must apply a torque, and $\Delta\vL$ points along the torque, not along $\vL$.
A fast-spinning top has a large $\vL$ along its axis; gravity's torque, trying to topple it, is horizontal
and \emph{perpendicular} to $\vL$. A perpendicular torque cannot lengthen or shorten $\vL$, only swivel
it, so instead of falling, the axis slowly sweeps around in a cone. This sideways sweep is
\textbf{precession}, and it is why a fast top stays up. Spin a wheel fast enough and ordinary torques
barely deflect its large $\vL$, so its axis holds a fixed direction: this is the gyroscope of aircraft and
spacecraft instruments.
\end{remarkbox}

\subsection{Spin plus orbital: how a system's angular momentum splits}
A system may move bodily from place to place and rotate about its own centre at the same time: the Earth
orbits the Sun while spinning on its axis, and a thrown spanner's centre arcs along a parabola while the
spanner tumbles. The total angular momentum separates cleanly into those two parts, provided the
splitting is done about the centre of mass.

\begin{derivbox}{The spin--orbital split}
Write each particle's position relative to the centre of mass, $\vr_i=\vr_{\text{cm}}+\vr_i'$, where
$\vr_{\text{cm}}$ is the position of the CM, as in \chref{ch:com} (so
$\vv_i=\vv_{\text{cm}}+\vv_i'$, where $\vv_{\text{cm}}$ is the CM's velocity and primes mark motion relative to the CM), and substitute into $\vL=\sum_i m_i(\vr_i\ex\vv_i)$. Multiplying out
gives four terms; the two ``cross'' terms contain $\sum_i m_i\vr_i'$ and $\sum_i m_i\vv_i'$, both \emph{zero}
by the definition of the centre of mass (positions relative to the CM average to nothing, and the
second sum is the time derivative of the first). What survives, with $\vc{P}=M\vv_{\text{cm}}$ the total linear momentum, is
\[
\vL=\underbrace{\vr_{\text{cm}}\ex\vc{P}}_{\text{orbital}}
+\underbrace{\sum_i \vr_i'\ex m_i\vv_i'}_{\text{spin}}.
\]
\end{derivbox}

\begin{keybox}{Orbital and spin angular momentum}
The \textbf{orbital} part, $\vr_{\text{cm}}\ex\vc{P}$, is the angular momentum of the whole mass
imagined concentrated at the CM; it tracks where the CM \emph{goes}. The \textbf{spin} part is the
angular momentum of the motion \emph{about} the CM; it tracks how the body turns on itself. The Earth's
total angular momentum about the Sun is its orbital part (once a year) plus its spin part (once a day);
each is conserved separately when the torque that would change it vanishes. This echoes \chref{ch:com}, where the CM carried the bodily motion of a system.
\end{keybox}

\subsection{Conservation of angular momentum in the wild}

With $\vtau_{\text{ext}}=d\vL/dt$ in hand, conservation follows at once: no external torque, no change in
$\vL$. For a body whose shape can change, that means $I_1\omega_1=I_2\omega_2$, and the skater is the
textbook case.

\begin{figure}[htbp]
\centering
\tikzfig{angmom_skater}
\figcaption{fig:skater}{A skater pulls their arms in and spins faster, with no one pushing them. Gravity and the normal force are both parallel to the vertical spin axis and the ice is nearly frictionless, so there is no external torque about that axis and $L=I\omega$ is fixed (the two $\vL$ arrows are equal). Pulling the arms in lowers $I$, so $\omega$ must rise to compensate.}
\end{figure}

\begin{trapbox}{angular momentum is conserved, kinetic energy is not}
The skater's kinetic energy \emph{rises} when she pulls in ($K=\tfrac{1}{2}I\omega^2=L^2/2I$, so halving
$I$ doubles $K$). It comes from her muscles: drawing the arms in needs an inward force acting through an
inward displacement, which does positive work. Conservation of $L$ says nothing about energy; the two are
independent laws.
\end{trapbox}

The same line, $I\omega$ before equals $I\omega$ after when no external torque acts, runs through a wide range of phenomena.

\begin{remarkbox}{Where conservation of angular momentum shows up}
\textbf{Planets.} Gravity pulls a planet straight towards the Sun, so it exerts no torque about the Sun,
and the planet's angular momentum about the Sun is conserved. That \emph{is} Kepler's second law
(\chref{ch:kepler}).\\[3pt]
\textbf{Collapsing stars.} A dying star's core, a few thousand kilometres across, collapses to a neutron
star about \SI{20}{km} across. With $I\propto R^2$ falling by $\sim10^4$ to $10^5$, $\omega$ rises by the
same factor: a core turning once in half an hour ends up spinning tens of times a \emph{second}.\\[3pt]
\textbf{The falling cat.} A dropped cat with zero angular momentum still turns feet-down by rotating its
front and rear halves in opposite senses, with different $I$; the total stays zero throughout.
\end{remarkbox}

\begin{trapbox}{``no external torque'' is about torque, not force}
Angular momentum is conserved when the net external \emph{torque} is zero, a condition independent of
whether the net force is zero. Gravity acts on the skater, the diver, and the falling cat the whole time,
but in a uniform field it acts \emph{at the centre of mass}, so its torque about the CM is zero, and $L$
about the CM is conserved regardless. Always check torque about your chosen axis, not whether forces are
present.
\end{trapbox}

\subsection{The crack: angular momentum is not continuous}

Everything above treats $\vL=I\vomega$ as a quantity that can take \emph{any} value. Spin the skater
slower and slower and $L$ slides smoothly towards zero, through every value in between. For an electron in
an atom, it does not. Angular momentum comes in indivisible \emph{lumps} of
$\hbar=h/2\pi=\SI{1.05e-34}{J.s}$, with $h$ Planck's constant: any measured component of an electron's
orbital angular momentum is a whole-number multiple of $\hbar$, never $1.3\hbar$, and its intrinsic
\emph{spin} has a component of $\pm\tfrac12\hbar$ along any axis. Nothing in this chapter can express such
a staircase of values. We never notice it because the step is tiny: a
spinning skater has $L\sim\SI{10}{kg.m^2/s}$, about $10^{35}$ steps, so at human scale $L$ looks
perfectly continuous, as a beach looks smooth until you pick up a handful of sand.

\why{Notice \emph{which half survives}. Conservation of $\vL$ holds \emph{exactly} in quantum mechanics. What fails is only the unspoken assumption that
$L$ is \emph{continuous}. The deep conservation law outlives the picture that produced it. In Bohr's model the quantisation of $L$ is the single new rule that, added to
\chref{ch:circdyn}'s circular dynamics (electrostatic attraction supplying the $mv^2/r$), gives the size
of the hydrogen atom and its spectrum.}

\section{Worked Examples}

The examples follow the order of the chapter: torque and equilibrium, moment of inertia and the
theorems, $\tau=I\alpha$, and finally angular momentum. Each names its principle before any arithmetic.

\begin{exbox}{Torque from components}
A force $\vF=3\ihat-\jhat$ (newtons) acts at the point $\vr=2\ihat+3\jhat$ (metres) from the origin. Find
the torque about the origin.

\stepn{1} \textbf{Set up the determinant.} Both vectors lie in the $xy$-plane, so only the $\khat$
component survives:
\[
\vtau=\vr\ex\vF=\begin{vmatrix}\ihat&\jhat&\khat\\ 2&3&0\\ 3&-1&0\end{vmatrix}
=\khat\big[(2)(-1)-(3)(3)\big].
\]

\stepn{2} \textbf{Evaluate.}
\[
\vtau=-11\,\khat\ \si{N.m}.
\]

\stepn{3} \textbf{Read the sign.} The negative $\khat$ means the torque points \emph{into} the page, so
it would turn the body \textbf{clockwise}. \ans{$\vtau=-11\khat$~\si{N.m}, clockwise}
\end{exbox}

\begin{exbox}{A plank on two supports}\label{ex:plank}
A uniform plank \SI{4}{m} long, of mass \SI{20}{kg}, rests on supports at its two ends $A$ and $B$. A
person of mass \SI{60}{kg} stands \SI{1}{m} from $A$ (Figure~\ref{fig:plank}). Find the force each
support exerts ($g=\SI{10}{m/s^2}$).

\stepn{1} \textbf{Principle.} The plank is in equilibrium, so the net force and the net torque about
\emph{any} point are both zero. The plank's weight, \SI{200}{N}, acts at its centre of mass, \SI{2}{m}
from $A$.

\stepn{2} \textbf{Take moments about $A$}, so that the unknown $N_A$ has no moment arm and drops out.
Anticlockwise (from $N_B$) equals clockwise (from the two weights):
\[
N_B(4)=600(1)+200(2)=1000\qquad\Longrightarrow\qquad N_B=\SI{250}{N}.
\]

\stepn{3} \textbf{Vertical forces balance:} $N_A=800-250=\SI{550}{N}$.

\stepn{4} \textbf{Check about $B$:} $N_A(4)=600(3)+200(2)=2200$, so $N_A=\SI{550}{N}$ \checkmark. The
support nearer the person carries more of the load, as experience suggests.
\ans{$N_A=\SI{550}{N}$, $N_B=\SI{250}{N}$}
\end{exbox}

\begin{figure}[htbp]
\centering
\tikzfig{torque_plank}
\figcaption{fig:plank}{A plank on two supports, arrows drawn to scale. Taking moments about $A$ removes
the unknown $N_A$ from the equation at once; the support nearer the person carries more than two thirds
of the total \SI{800}{N}.}
\end{figure}

\begin{exbox}{Four masses at the corners of a square}
Four particles, each of mass \SI{2}{kg}, sit at the corners of a square of side $a=\SI{2}{m}$. Find the
moment of inertia about (a) an axis through the centre perpendicular to the plane, and (b) an axis along
one side.

\stepn{1} \textbf{(a) Central perpendicular axis.} Every mass is at the same distance, namely half the
diagonal:
\[
r=\frac{a\sqrt2}{2}=\frac{2\sqrt2}{2}=\sqrt2\ \si{m}.
\]
Hence
\[
I=\sum mr^{2}=4(2)(\sqrt2)^{2}=4(2)(2)=\SI{16}{kg.m^2}.
\]

\stepn{2} \textbf{(b) Axis along one side.} Now two masses lie \emph{on} the axis and contribute
nothing, while the other two are at the full side length $a=\SI{2}{m}$:
\[
I=2(2)(2)^{2}=\SI{16}{kg.m^2}.
\]

\stepn{3} \textbf{Comment.} The two answers coincide here by arithmetic accident; it is not a general
rule. Always locate each mass's distance \emph{from the stated axis} before summing.
\ans{\SI{16}{kg.m^2} in both cases}
\end{exbox}

\begin{exbox}{Choosing the right theorem}
Find the moment of inertia of a uniform disc of mass $M$ and radius $R$ about (a) a tangent lying
\emph{in} its plane, and (b) a tangent \emph{perpendicular} to its plane.

\stepn{1} \textbf{(a) Tangent in the plane.} This is a shift from a \emph{diameter}, so we need the
diameter value first. By the perpendicular-axis theorem on the flat disc,
$I_z=I_x+I_y$ with $I_z=\tfrac{1}{2}MR^{2}$ and $I_x=I_y$ by symmetry, giving
$I_{\text{diameter}}=\tfrac{1}{4}MR^{2}$. Now shift by $d=R$ using the parallel-axis theorem:
\[
I=\tfrac{1}{4}MR^{2}+MR^{2}=\tfrac{5}{4}MR^{2}.
\]

\stepn{2} \textbf{(b) Tangent perpendicular to the plane.} Start from the central perpendicular axis,
$I_{\text{cm}}=\tfrac{1}{2}MR^{2}$, and shift by $d=R$:
\[
I=\tfrac{1}{2}MR^{2}+MR^{2}=\tfrac{3}{2}MR^{2}.
\]

\stepn{3} \textbf{The lesson.} Both parts used the parallel-axis theorem, but part (a) needed the
perpendicular-axis theorem \emph{first} to obtain a starting value. Identify the axis through the centre
of mass that is \emph{parallel} to your target axis before you shift.
\ans{$\tfrac{5}{4}MR^{2}$ and $\tfrac{3}{2}MR^{2}$}
\end{exbox}

\begin{exbox}{A flywheel under constant torque}
A torque of \SI{20}{N.m} acts on a flywheel of moment of inertia \SI{5}{kg.m^2}, initially at rest.
Find (a) the angular acceleration, (b) the angular speed after \SI{6}{s}, and (c) the angle turned in
that time.

\stepn{1} \textbf{(a)} From $\tau=I\alpha$:
\[
\alpha=\frac{\tau}{I}=\frac{20}{5}=\SI{4}{rad/s^2}.
\]

\stepn{2} \textbf{(b)} With constant $\alpha$ from rest,
$\omega=\alpha t=4(6)=\SI{24}{rad/s}$.

\stepn{3} \textbf{(c)} $\theta=\tfrac{1}{2}\alpha t^{2}=\tfrac{1}{2}(4)(36)=\SI{72}{rad}$, which is
$72/2\pi\approx11.5$~revolutions.

\stepn{4} \textbf{Check by energy.} The work done is $\tau\theta=20(72)=\SI{1440}{J}$, and the kinetic
energy gained is $\tfrac{1}{2}I\omega^{2}=\tfrac{1}{2}(5)(576)=\SI{1440}{J}$ \checkmark
\ans{\SI{4}{rad/s^2}; \SI{24}{rad/s}; \SI{72}{rad}}
\end{exbox}

\begin{exbox}{A mass unwinding a cord from a wheel}
A uniform disc of mass \SI{10}{kg} and radius \SI{0.2}{m} is free to rotate about a fixed horizontal
axis. A light cord wound round its rim carries a hanging mass of \SI{2}{kg}. Find the acceleration of
the hanging mass and the tension in the cord ($g=\SI{10}{m/s^2}$).

\stepn{1} \textbf{Moment of inertia of the disc.}
\[
I=\tfrac{1}{2}MR^{2}=\tfrac{1}{2}(10)(0.2)^{2}=\SI{0.2}{kg.m^2}.
\]

\stepn{2} \textbf{Two equations, one for each body.} For the hanging mass (taking down as positive):
\[
mg-T=ma. \tag{i}
\]
For the disc, the cord's tension exerts a torque $TR$. Because the cord does not slip on the rim, the rim's
tangential acceleration equals the cord's, $a=\alpha R$, so $\alpha=a/R$:
\[
TR=I\alpha=I\frac{a}{R}
\qquad\Longrightarrow\qquad
T=\frac{Ia}{R^{2}}. \tag{ii}
\]

\stepn{3} \textbf{Eliminate $T$.} Substituting (ii) into (i):
\[
mg=ma+\frac{Ia}{R^{2}}
\qquad\Longrightarrow\qquad
a=\frac{mg}{m+I/R^{2}}=\frac{2(10)}{2+0.2/0.04}=\frac{20}{7}=\SI{2.86}{m/s^2}.
\]

\stepn{4} \textbf{Tension.}
\[
T=m(g-a)=2(10-2.857)=\SI{14.3}{N}.
\]

\stepn{5} \textbf{Cross-check} with (ii): $T=\dfrac{0.2(2.857)}{0.04}=\SI{14.3}{N}$ \checkmark. Note
$T<mg=\SI{20}{N}$, as it must be for a descending mass, and $a<g$ because the disc's inertia resists.
\ans{$a=\SI{2.86}{m/s^2}$; $T=\SI{14.3}{N}$}
\end{exbox}

\begin{exbox}{The skater, quantitatively}
A skater spinning at \SI{2}{rev/s} with moment of inertia \SI{6}{kg.m^2} pulls her arms in, reducing it
to \SI{2}{kg.m^2}. Find her new rate of spin and the factor by which her kinetic energy changes. Where
does the extra energy come from?

\stepn{1} \textbf{No external torque} about the vertical axis (friction at the skate is negligible), so
$I\omega$ is conserved:
\[
I_1\omega_1=I_2\omega_2
\qquad\Longrightarrow\qquad
\omega_2=\frac{6(2)}{2}=\SI{6}{rev/s}.
\]
She spins three times faster.

\stepn{2} \textbf{Kinetic energy.} Writing $K=\dfrac{L^{2}}{2I}$ with $L$ fixed, $K\propto1/I$. Since
$I$ falls to one third,
\[
\frac{K_2}{K_1}=\frac{I_1}{I_2}=3 .
\]

\stepn{3} \textbf{Where the energy came from.} From her muscles, as the caution after Figure~\ref{fig:skater}
explained: pulling the arms in does positive work. \ans{\SI{6}{rev/s}; $K$ triples}
\end{exbox}

\begin{exbox}{A man walking on a turntable}
A man of mass \SI{50}{kg} stands on the rim of a turntable of moment of inertia \SI{200}{kg.m^2} and
radius \SI{1}{m}, both initially at rest on frictionless bearings. The man begins to walk along the rim
at \SI{1}{m/s} relative to the ground. Find the angular speed of the turntable.

\stepn{1} \textbf{Identify the conserved quantity.} There is no external torque about the vertical
axis, and the system starts at rest, so the \emph{total} angular momentum stays zero throughout.

\stepn{2} \textbf{Write the two contributions.} The man, treated as a particle on the rim, carries
$L_{\text{man}}=mvR$. The turntable carries $L_{\text{table}}=I\omega$. Their sum must vanish:
\[
I\omega+mvR=0 .
\]

\stepn{3} \textbf{Solve.}
\[
\omega=-\frac{mvR}{I}=-\frac{50(1)(1)}{200}=\SI{-0.25}{rad/s}.
\]
The minus sign says the turntable rotates \textbf{opposite} to the man's walk, which is the whole point
of the problem. \ans{\SI{0.25}{rad/s}, opposite to the man}
\end{exbox}

\begin{exbox}{A hinged rod falling from horizontal}
A uniform rod of mass $M$ and length $\ell$ is hinged at one end and released from rest in a horizontal
position. Find its angular speed as it passes through the vertical, and evaluate for
$\ell=\SI{1.2}{m}$ ($g=\SI{10}{m/s^2}$).

\stepn{1} \textbf{Use energy, not torque.} The torque on the rod changes continuously as it falls, so
$\alpha$ is not constant and the constant-$\alpha$ equations do not apply; $\tau=I\alpha$ holds at each
instant but would have to be integrated. Energy conservation bypasses that entirely.

\stepn{2} \textbf{The drop of the centre of mass.} The centre of mass is at $\ell/2$ from the hinge, so
in swinging from horizontal to vertical it descends $h=\ell/2$:
\[
\text{PE lost}=Mg\frac{\ell}{2}.
\]

\stepn{3} \textbf{The kinetic energy gained} is wholly rotational about the hinge, with
$I=\tfrac{1}{3}M\ell^{2}$:
\[
Mg\frac{\ell}{2}=\tfrac{1}{2} I\omega^{2}=\tfrac{1}{2}\left(\tfrac{1}{3}M\ell^{2}\right)\omega^{2}.
\]

\stepn{4} \textbf{Solve.} The mass cancels:
\[
\omega=\sqrt{\frac{3g}{\ell}}.
\]
For $\ell=\SI{1.2}{m}$: $\omega=\sqrt{30/1.2}=\sqrt{25}=\SI{5}{rad/s}$.

\stepn{5} \textbf{A caution.} The tip of the rod then moves at $v=\omega\ell=\SI{6}{m/s}$, which is
\emph{faster} than a particle dropped through the tip's own fall of $\ell$
($\sqrt{2g(1.2)}=\SI{4.90}{m/s}$). The rigid rod carries energy released by its inner parts out to the
tip, so the tip gets more than its own share. \ans{$\omega=\sqrt{3g/\ell}=\SI{5}{rad/s}$}
\end{exbox}

\section{Summary}\label{sec:summary14}
Rotation has its own versions of position, velocity, acceleration, mass, force and momentum, and they obey
equations of exactly the same form as their translational partners. The only genuinely new object is the
moment of inertia $I$, and unlike mass it depends on \emph{which axis} is chosen, so the axis must be named
before anything is written down. The table below collects the results.

\medskip
{\small\renewcommand{\arraystretch}{1.45}
\begin{center}
\begin{tabular}{@{}ll@{}}
\toprule
Quantity & Result\\
\midrule
Torque             & $\vtau=\vr\ex\vF$, \quad $\tau=rF\sin\theta=F r_\perp$\\[2pt]
Moment of inertia  & $I=\sum m_ir_i^{2}=\int r^{2}\,dm$\\[2pt]
Second law         & $\vtau_{\text{ext}}=\dfrac{d\vL}{dt}$; \quad $\tau=I\alpha$ about a fixed axis\\[4pt]
Angular momentum   & $\vL=\vr\ex\vp$; \quad $L=I\omega$ for a rigid body\\[2pt]
Rotational KE      & $K=\tfrac{1}{2}I\omega^{2}$\\[2pt]
Work and power     & $W=\int\tau\,d\theta$ ($=\tau\theta$ for constant $\tau$), \quad $P=\tau\omega$\\[2pt]
Parallel axis      & $I=I_{\text{cm}}+Md^{2}$ \quad (\emph{any} body)\\[2pt]
Perpendicular axis & $I_z=I_x+I_y$ \quad (\emph{laminae only})\\[2pt]
Conservation       & $\vtau_{\text{ext}}=\vc{0}\Rightarrow I_1\omega_1=I_2\omega_2$\\
\bottomrule
\end{tabular}
\end{center}
}

\medskip
\textbf{The standard moments} ($I=cMR^{2}$, or $cM\ell^{2}$ for a rod): ring $1$; disc or solid
cylinder $\tfrac{1}{2}$; solid sphere $\tfrac{2}{5}$; spherical shell $\tfrac{2}{3}$; rod about centre
$\tfrac{1}{12}$; rod about one end $\tfrac{1}{3}$; disc about a diameter $\tfrac{1}{4}$.

\medskip
\textbf{Three questions that settle most problems.} \emph{(1)} About which axis?
\emph{(2)} Is there an external \emph{torque}? A non-zero net force can still give zero torque.
\emph{(3)} Does the body also translate? Then split $\vL$ into orbital plus spin about the centre of
mass.

\section{Problem Types}\label{sec:types14}
Six problem types recur in rotational mechanics. The first is the general case; the others are its most
common forms.

\archq{Rotational-motion problems}
{A rotating rigid body; find $I$, angular acceleration, or a final angular speed.}
{Name the axis first ($I$ is meaningless without it). Look up $I$, shift with parallel-axis if needed.
For dynamics use $\tau=I\alpha$; for a final speed after a shape change or drop, use $L$ conservation or
energy $K=\tfrac{1}{2}I\omega^2$. A force through the axis (including gravity at the CM) gives no torque.}
{Using the constant-$\alpha$ equations when $\tau$ (hence $\alpha$) changes through the motion, as for a
swinging rod: $\tau=I\alpha$ still holds at each instant, but use energy for final speeds.}
{A uniform disc of mass $M$ and radius $R$, free to turn about its central axis, is acted on by a torque
$\tau$ about that axis. Its angular acceleration is:}
{$\tau/MR^2$}{$\tau/2MR^2$}{$2\tau/3MR^2$}{$2\tau/MR^2$\archans{D}{About the central axis $I=\tfrac12MR^2$, so
$\alpha=\tau/I=2\tau/MR^2$. (A) uses the ring's $I=MR^2$; (B) multiplies by the factor $\tfrac12$ instead of dividing by
it; (C) uses $I=\tfrac32MR^2$, the moment of inertia about a parallel axis through the rim, not the centre.}}

\archq{Angular kinematics at constant angular acceleration}
{Constant angular acceleration: given some of $\omega_0,\omega,\alpha,\theta,t$, find another. The
rotational twin of \chref{ch:straightline}.}
{Use the exact analogues of the straight-line equations: $\omega=\omega_0+\alpha t$,
$\theta=\omega_0 t+\tfrac{1}{2}\alpha t^2$, and $\omega^2=\omega_0^2+2\alpha\theta$. If $\alpha$ varies with
time, integrate instead ($\omega=\int\alpha\,dt$, $\theta=\int\omega\,dt$).}
{Mixing revolutions with radians ($1$ rev $=2\pi$ rad), or using the constant-$\alpha$ formulas when
$\alpha$ is not constant.}
{A wheel starting from rest reaches \SI{20}{rad/s} in \SI{4}{s} under constant angular acceleration. How
many radians does it turn in that time?}
{\SI{160}{rad}}{\SI{80}{rad}}{\SI{40}{rad}}{\SI{20}{rad}\archans{C}{$\alpha=20/4=\SI{5}{rad/s^2}$, then
$\theta=\tfrac{1}{2}\alpha t^2=\tfrac{1}{2}(5)(16)=\SI{40}{rad}$ (or use $\theta=\bar\omega t=10\times4$).}}

\archq{One body dropped onto another, angular-momentum conservation}
{A rotating body has a second body dropped onto it (a disc on a disc, putty on a turntable, a person
stepping on a merry-go-round); find the common final angular speed.}
{No external torque about the axis, so $\vL$ is conserved: $I_1\omega_1=(I_1+I_2)\omega_2$. Add the
dropped body's moment of inertia $I_2$ (about the same axis) and solve. Kinetic energy is \emph{not}
conserved: the coupling is an inelastic ``rotational collision''.}
{Assuming kinetic energy is conserved (it is dissipated as heat while the surfaces slip up to a common speed), or forgetting
to put $I_2$ about the \emph{shared} axis.}
{A disc ($I=\tfrac{1}{2}MR^2$) spins at $\omega$. An identical disc, at rest, is dropped coaxially onto it.
The common angular speed is:}
{$\omega$}{$\dfrac{\omega}{2}$}{$\dfrac{2\omega}{3}$}{$\dfrac{\omega}{4}$\archans{B}{Both discs have the
same $I$, so $I\omega=(I+I)\omega_2\Rightarrow\omega_2=\omega/2$. Half the kinetic energy is lost to
friction between them. (If the second disc had the same radius but mass $M/2$, the answer would instead be $\tfrac{2}{3}\omega$.)}}

\archq{Something strikes a pivoted rod: angular momentum about the pivot}
{A particle, bullet, or lump hits a rod (or door, or rod-on-a-table) that can rotate about a pivot, and
embeds or bounces; find the angular speed just after.}
{Conserve angular momentum \emph{about the pivot} (the pivot's force has zero moment arm there). Before:
$L=mv\,d$ with $d$ the perpendicular distance of the line of motion from the pivot. After:
$L=I_{\text{total}}\omega$, with $I_{\text{total}}$ the rod's $I$ about the pivot plus $md^2$ for an
embedded mass. Do \emph{not} conserve linear momentum: the pivot supplies an external force.}
{Using linear-momentum conservation (the pivot breaks it), or taking angular momentum about the wrong
point.}
{A bullet of mass $m$, speed $v$, embeds in the tip of a uniform rod (mass $M$, length $\ell$) pivoted at
its other end. The rod's angular speed just after is:}
{$\dfrac{3mv}{(M+3m)\ell}$}{$\dfrac{mv}{M\ell}$}{$\dfrac{2mv}{(M+2m)\ell}$}{$\dfrac{mv}{(M+m)\ell}$\archans{A}{Angular
momentum about the pivot: $mv\ell=\big(\tfrac{1}{3}M\ell^2+m\ell^2\big)\omega$, so
$\omega=\dfrac{mv\ell}{(\tfrac{1}{3}M+m)\ell^2}=\dfrac{3mv}{(M+3m)\ell}$. (B) treats the rod as a point mass at
its tip and forgets the bullet; (C) uses $\tfrac12M\ell^2$ for the rod; (D) uses $M\ell^2$.}}

\archq{A hanging mass unwinds a cord from a wheel}
{A mass on a string wrapped around a pulley, disc, or cylinder that has its own moment of inertia; find
the acceleration, tension, or angular speed after a drop.}
{The pulley's inertia matters, so it is \emph{two} equations: $mg-T=ma$ for the falling mass and
$TR=I\alpha=I(a/R)$ for the wheel, linked by $a=\alpha R$. Eliminate $T$. For a disc pulley this gives
$a=\dfrac{g}{1+I/mR^2}=\dfrac{mg}{m+M/2}$.}
{Ignoring the pulley's moment of inertia. A massless pulley would give $T=0$ and $a=g$; a real pulley makes $0<T<mg$ and $a<g$.}
{A mass $m$ hangs from a cord wrapped around a disc pulley of mass $M$. Its downward acceleration is:}
{$g$}{$\dfrac{mg}{m+M}$}{$\dfrac{mg}{m+M/2}$}{$\dfrac{mg}{m+2M}$\archans{C}{$mg-T=ma$ and
$T=\tfrac{1}{2}Ma$ (from $TR=\tfrac{1}{2}MR^2\cdot a/R$); adding, $mg=(m+\tfrac{1}{2}M)a$, so
$a=\dfrac{mg}{m+M/2}$. With $M=0$ this correctly returns $a=g$.}}

\archq{Moment of inertia about a shifted or in-plane axis}
{``Find $I$ of a standard body about a tangent, an edge, a diameter, or an axis through the end.''}
{Start from the standard $I_{\text{cm}}$, then apply the right theorem. \textbf{Parallel axis}
($I=I_{\text{cm}}+Md^2$) works for \emph{any} body and shifts to a parallel axis a distance $d$ away.
\textbf{Perpendicular axis} ($I_z=I_x+I_y$) works only for \emph{flat, laminar} bodies.}
{Applying the perpendicular-axis theorem to a sphere or a cylinder; it is valid for laminae only.
Also, $d$ is measured from the axis through the \emph{centre of mass}, not from any convenient point.}
{The moment of inertia of a solid sphere of mass $M$ and radius $R$ about a tangent is:}
{$\tfrac{2}{5}MR^2$}{$\tfrac{7}{5}MR^2$}{$\tfrac{3}{2}MR^2$}{$\tfrac{2}{3}MR^2$\archans{B}{Parallel axis with
$I_{\text{cm}}=\tfrac{2}{5}MR^2$ and $d=R$: $I=\tfrac{2}{5}MR^2+MR^2=\tfrac{7}{5}MR^2$.}}

\section{Practice Problems}\label{sec:practice14}

The problems below are of standard examination difficulty and are meant to be worked without looking back at the chapter. Answers, with brief working, follow the list.

\begin{enumerate}[itemsep=3pt]
\item Find the moment of inertia of a uniform rod of mass $M$ and length $\ell$ about a perpendicular
      axis through a point $\ell/4$ from its centre.
\item Find the moment of inertia of a thin uniform ring of mass $M$, radius $R$, about a tangent in its
      plane.
\item A torque of \SI{10}{N.m} acts on a flywheel of moment of inertia \SI{2}{kg.m^2} at rest. Find
      the angular acceleration and the time to reach \SI{50}{rad/s}.
\item Masses \SI{3}{kg} and \SI{5}{kg} are fixed at the ends of a light rod of length \SI{2}{m}. Find
      the moment of inertia about a perpendicular axis through the centre of mass.
\item A skater spinning at \SI{1}{rev/s} with $I=\SI{4}{kg.m^2}$ pulls in to $I=\SI{1}{kg.m^2}$. Find
      the new rate of spin and the factor by which the kinetic energy changes.
\item A uniform disc of mass \SI{2}{kg} and radius \SI{0.5}{m} spins about its central axis at
      \SI{10}{rad/s}. Find its angular momentum.
\end{enumerate}

\subsection*{Answers}
{\small
\begin{enumerate}[itemsep=2pt]
\item $\tfrac{1}{12}M\ell^{2}+M(\ell/4)^{2}
=\left(\tfrac{1}{12}+\tfrac{1}{16}\right)M\ell^{2}=\tfrac{7}{48}M\ell^{2}$.
\item Perpendicular axis first ($I_z=MR^{2}=2I_{\text{diameter}}$) gives $I_{\text{diameter}}=\tfrac{1}{2}MR^{2}$;
then parallel axis with $d=R$: $\tfrac{1}{2}MR^{2}+MR^{2}=\tfrac{3}{2}MR^{2}$.
\item $\alpha=\tau/I=\SI{5}{rad/s^2}$; $t=\omega/\alpha=\SI{10}{s}$.
\item The centre of mass lies \SI{1.25}{m} from the \SI{3}{kg} and \SI{0.75}{m} from the
\SI{5}{kg}, so $I=3(1.25)^{2}+5(0.75)^{2}=\SI{7.5}{kg.m^2}$.
\item $I\omega$ fixed gives \SI{4}{rev/s}; and $K=L^{2}/2I$ with $L$ fixed, so $K$ rises by a
factor of $4$, the skater supplying that energy by muscular work.
\item $I=\tfrac{1}{2}MR^{2}=\SI{0.25}{kg.m^2}$, so $L=I\omega=\SI{2.5}{kg.m^2/s}$.
\end{enumerate}
}
